% !PRINTED 36
\weberpage{36}
Wenn $a$ und $b$ weiter in Factoren zerlegbar sind, die zu einander relativ prim sind, so können wir mit den Gruppen $\mathfrak{A}$ und $\mathfrak{B}$ wieder ebenso verfahren, und wir kommen also zu dem Resultate, dass unser Theorem I. allgemein bewiesen ist, wenn wir es noch für Gruppen nachweisen können, deren Grad eine Potenz einer Primzahl ist.

Es sei also jetzt der Grad $n$ der Gruppe $S$ eine Potenz einer Primzahl $p$
\weq{9}{n = p^k.}
Die Grade aller Elemente von $S$, die ja Divisoren von $n$ sind, müssen dann gleichfalls Potenzen von $p$ sein. Es ist nachzuweisen, dass eine solche Gruppe durch eine Basis darstellbar ist.

Man wähle in $S$ ein Element $A$ von möglichst hohem Grade $a$. Dann ist $a$ eine Potenz von $p$ und die Grade aller anderen Elemente sind Theiler von $a$, so dass für jedes Element $\Theta$ von $S$
\weq{10}{\Theta^a = 1}
ist. Die Elemente
\weq{11}{1, A, A^2 \ldots A^{a-1}}
sind alle von einander verschieden, und ihre Gesammtheit ist ein Theiler von $S$. Ist damit die Gruppe $S$ erschöpft, ist also jedes Element von der Form $A^\alpha$, so ist $S$ durch eine eingliedrige Basis dargestellt. Wenn aber $S$ mit (11) noch nicht erschöpft ist, so wird es doch für jedes Element $\Theta$ von $S$ einen gewissen Exponenten $h$ geben, so dass $\Theta^h$ in der Reihe (11) enthalten ist. Gewiss wird das eintreten, wenn $h$ der Grad von $\Theta$, also $\Theta^h = 1$ ist. Unter diesen Zahlen $h$ wird eine die kleinste positive sein, die wir mit $b$ bezeichnen wollen. Es giebt also für jedes Element $\Theta$ eine gewisse kleinste positive Zahl $b$, so dass
\weq{12}{\Theta^b = A^\lambda}
in der Reihe (11) enthalten ist.

Diese Zahl $b$ ist ein Theiler von $a$, also auch eine Potenz von $p$ und zugleich ein Theiler von $\lambda$. Denn setzen wir $a = qb + b'$, wo $0 \leqq b' < b$ ist, so folgt aus (10) und (12)
\[
\Theta^a = A^{\lambda q} \Theta^{b'} = 1,
\]
oder $\Theta^{b'} = A^{-\lambda q}$, und wenn also $b'$ nicht Null ist, so giebt es gegen die Voraussetzung eine noch kleinere Zahl als $b$, nämlich $b'$, die der Forderung (12) genügt.
